\section*{Paraguai em Perspectiva Regional}\label{comparativo-paises}
\addcontentsline{toc}{section}{Paraguai em Perspectiva Regional}

Para situar o Paraguai no contexto de seus vizinhos imediatos, a tabela abaixo compara os principais critérios relevantes para relocação estratégica. Os dados referem-se a 2023--2024 e são de fontes públicas (Banco Mundial, CEPAL, UNODC, KPMG, Mercer Cost of Living).

\vspace{0.5em}

\begin{center}
\resizebox{\linewidth}{!}{%
{\renewcommand{\arraystretch}{1.35}
\begin{tabular}{@{}l c c c c c c@{}}
\toprule
\textbf{Critério} & \textbf{Paraguai} & \textbf{Brasil} & \textbf{Argentina} & \textbf{Bolívia} & \textbf{Chile} & \textbf{Uruguai} \\
\midrule
Carga tributária (\% PIB) & \textbf{15\%} & 33\% & 29\% & 23\% & 21\% & 28\% \\
Imposto de renda PF (teto) & \textbf{10\%} & 27,5\% & 35\% & 25\% & 40\% & 36\% \\
ISS/IVA (consumo) & \textbf{10\%} & 12--25\% & 21\% & 13\% & 19\% & 22\% \\
\midrule
Custo da terra agrícola (USD/ha) & \textbf{1.500--5.000} & 8.000--25.000 & 3.000--15.000 & 800--3.000 & 8.000--30.000 & 5.000--12.000 \\
Acesso à propriedade rural para estrangeiros & \textbf{Amplo} & Restrito & Moderado & Moderado & Muito restrito & Moderado \\
\midrule
Homicídios (por 100 mil hab.) & \textbf{5,3} & 22,3 & 5,0 & 6,1 & 4,5 & 10,8 \\
Risco geopolítico regional & Baixo & Médio & Médio & Médio & Baixo & Muito baixo \\
Estabilidade da moeda & \textbf{Alta} & Baixa & Muito baixa & Moderada & Alta & Alta \\
\midrule
Custo de vida relativo (PY = 100) & \textbf{100} & 165 & 120 & 75 & 205 & 175 \\
Energia elétrica (USD/kWh) & \textbf{0,07} & 0,12 & 0,06$^{\dagger}$ & 0,08 & 0,18 & 0,18 \\
\midrule
Facilidade de residência permanente & \textbf{Alta} & Moderada & Moderada & Moderada & Baixa & Alta \\
Processo de naturalização (anos mínimos) & \textbf{3} & 4 & 2 & 2 & 5 & 5 \\
Idioma (barreira para anglófonos) & Espanhol/\allowbreak Guarani & Português & Espanhol & Espanhol & Espanhol & Espanhol \\
\midrule
Conectividade (% pop. com internet) & 68\% & 84\% & 86\% & 48\% & 92\% & 88\% \\
Cobertura 4G rural & Moderada & Moderada & Moderada & Baixa & Alta & Alta \\
\bottomrule
\end{tabular}}%
}
\end{center}

\vspace{0.3em}
{\small $^{\dagger}$ A Argentina subsidia fortemente a energia elétrica, tornando o custo nominal baixo, mas a instabilidade da rede e os apagões frequentes são fatores negativos.}

\vspace{0.8em}

\textbf{Leitura do quadro comparativo}

O Paraguai destaca-se em três dimensões objetivamente mensuráveis: \textbf{carga tributária} (a mais baixa da América do Sul para pessoas físicas e jurídicas), \textbf{acesso e custo da terra} (o mais acessível entre os países com clima e solo favoráveis) e \textbf{facilidade de residência legal} (processo relativamente simples e rápido para estrangeiros). A criminalidade se situa em nível moderado, abaixo da média regional, com exceção de zonas de fronteira específicas detalhadas neste guia.

As principais limitações comparativas são: conectividade rural ainda em desenvolvimento, serviços de saúde pública de qualidade heterogênea, e dependência logística de infraestrutura ainda em expansão no interior. Quem busca ambiente urbano de padrão europeu encontrará melhor oferta no Uruguai ou Chile — pagando preço proporcional em custo de vida e tributação.

\medskip
O Paraguai representa o ponto de equilíbrio entre \textbf{acesso econômico}, \textbf{estabilidade institucional relativa} e \textbf{recursos naturais abundantes} para quem busca autossuficiência e qualidade de vida fora dos grandes centros urbanos.
